\subsection{Commutative stellar algebras}

Lemma 8. --- Let X and Y be metric spaces, with X complete. Let f be a mapping from X into Y such that

\[
d (f (x), f (y)) \geqslant d (x, y)
\]

for all x and y in X and such that the graph of f is closed in X × Y. Then f is a closed map.

Let F be a closed subset of X and let \((y_{n})_{n\in\mathbf{N}}\) a sequence in \(f(\mathrm{F})\) which converges to \(y\in Y\) . For every \(n\in N\) , let \(x_{n}\in F\) such that \(f(x_{n})=y_{n}\) . The hypothesis implies that the sequence \((x_{n})_{n\in\mathbf{N}}\) is a Cauchy sequence; let \(x\in X\) be its limit; it is an element of F, because F is closed. Moreover, we have \((x_{n},f(x_{n}))\to(x,y)\) into \(X\times Y\) . Since the graph of f is closed, it follows that \(y=f(x)\) belongs to \(f(\mathrm{F})\) .

Lemma 9. --- Let A be a commutative stellar algebra. Every character of A is Hermitian, and the Gelfand transformation is a morphism of stellar algebras from A into \(\mathcal{C}_{0}(\mathsf{X}(\mathsf{A}))\) .

It suffices to prove the first assertion (proposition 7 of I, p.~38 and lemma 3 of I, p.~98). Let \(\chi\) be a character of A. For every Hermitian element \(y\) of A, we have \(\chi(y) = \mathcal{G}_{\mathrm{A}}(y)(\chi) \in \mathrm{Sp}'(y) \subset \mathbf{R}\) (prop. 4 of I, p.~106). Consequently, the character \(\chi\) is Hermitian (prop. 1 of I, p.~95).

THEOREM 1. --- Let A be a commutative stellar algebra. The Gelfand transformation is an isometric isomorphism of the stellar algebra A onto the stellar algebra \(\mathcal{C}_{0}(\mathsf{X}(\mathsf{A}))\) of continuous functions on \(\mathsf{X}(\mathsf{A})\) vanishing at infinity.

The Gelfand transform is an involutive algebra morphism from A into \(\mathcal{C}_{0}(\mathsf{X}(\mathsf{A}))\) (lemma 9). Let B be its image. It is an involutive subalgebra of \(\mathcal{C}_{0}(\mathsf{X}(\mathsf{A}))\) . The elements of B separate the points of \(\mathsf{X}(\mathsf{A})\) by definition. Let \(\chi \in \mathsf{X}(\mathsf{A})\) . There exists \(x \in \mathsf{A}\) such that \(\chi(x) \neq 0\) , hence \(f = \mathcal{G}(x)\) is an element of B such that \(f(\chi) \neq 0\) . By TG, X, p.~40, cor. 2, the subalgebra B is therefore dense in \(\mathcal{C}_{0}(\mathsf{X}(\mathsf{A}))\) .

For every Hermitian element \(y\) of A, we have \(\| \mathcal{G}(y)\| = \varrho (y) = \| y\|\) (prop. 7 of I, p.~38 and formula (2) of I, p.~104), whence, for every \(x\in \mathrm{A}\) , the equalities

\[
\| x \| ^ {2} = \| x ^ {*} x \| = \| \mathcal {G} (x ^ {*} x) \| = \| \overline {{\mathcal {G} (x)}} \cdot \mathcal {G} (x) \| = \| \mathcal {G} (x) \| ^ {2}.
\]

Thus G is isometric, and its image B is therefore closed (Lemma 8). We conclude that \(\mathrm{B} = \mathcal{C}_{0}(\mathrm{X}(\mathrm{A}))\) .

COROLLARY 1. --- Let A be a stellar algebra and let x be a normal element of A. Then \(\|x\| = \varrho(x)\) .

Since the stellar subalgebra of A generated by x and \(x^{*}\) is commutative, we may assume that A is commutative. In this case, the corollary follows from Th. 1, Example 2 of I, p.~102, and Th. 1 of I, p.~24.

COROLLARY 2. --- Let A be a commutative stellar algebra.

\begin{enumerate}
\def\labelenumi{\alph{enumi})}
\item
  There exists a locally compact topological space X such that A is isomorphic to the stellar algebra \(\mathcal{C}_{0}(\mathrm{X})\) ;
\item
  Let B be a commutative stellar algebra. The map \(\pi \mapsto \mathsf{X}'(\pi)\) is a bijection from the set of stellar algebra morphisms from A to B onto the set of proper partial maps from \(\mathsf{X}(\mathsf{B})\) into \(\mathsf{X}(\mathsf{A})\) (def. 1 of I, p.~33).
\end{enumerate}

Theorem 1 establishes the first assertion, and the second assertion follows from Prop. 3 of I, p.~34 and from Example 2 of I, p.~102.

COROLLARY 3. --- Let A be a unital commutative stellar algebra.

\begin{enumerate}
\def\labelenumi{\alph{enumi})}
\item
  There exists a compact topological space X such that A is isomorphic to the stellar algebra \(\mathcal{C}(\mathrm{X})\) ;
\item
  Let B be a commutative unital stellar algebra. The map \(\pi \mapsto \mathsf{X}(\pi)\) is a bijection from the set of unital morphisms of stellar algebras from A to B to the set of continuous maps from \(\mathsf{X}(\mathsf{B})\) into \(\mathsf{X}(\mathsf{A})\) .
\end{enumerate}

This follows from the foregoing and from prop. 4 of I, p.~35.

Remark. --- *Let G be the category whose objects are locally compact spaces and whose morphisms are proper partial maps (def. 1 of I, p.~33), and let S be the category of commutative stellar algebras, whose morphisms are morphisms of involutive algebras. Consider the functor from S to the opposite category G° which associates to a commutative stellar algebra A the locally compact space X(A) of nonzero characters of A, and to a morphism φ: A → B of commutative stellar algebras the continuous map X'(φ). Th. 1 and Cor. 2 mean that this functor is an equivalence of categories, and that a quasi-inverse of this functor is the functor which associates to a locally compact topological space X the commutative stellar algebra C₀(X).

Similarly, Corollary 3 means that the opposite category of the category of compact spaces is equivalent to the category of unital commutative stellar algebras.*

